% !PRINTED 275
\weberpage{275}
Es lassen sich dann jeder Transformation $n$ten Grades von
der Form
\weq{2}{
\begin{pmatrix} a, & 0\\ c, & \partial \end{pmatrix}
}
je eine und nur eine Transformation der Grade $n'$, $n''$:
\weq{3}{
\begin{pmatrix} a', & 0\\ c', & \partial' \end{pmatrix}, \quad
\begin{pmatrix} a'', & 0\\ c'', & \partial'' \end{pmatrix}
}
zuordnen, die durch folgende Bedingungen bestimmt sind:
\weq{4}{
\begin{aligned}
a &= a'a'', & \partial &= \partial'\partial'',\\
\partial''c' &\equiv c \pmod{a'}, & \partial'c'' &\equiv c \pmod{a''},
\end{aligned}
}
und umgekehrt folgt aus jedem Paar Transformationen von der
Form (3) nach (4) eine und nur eine Transformation (2). Nach
(4) sind nämlich zunächst $\partial'$, $\partial''$ bestimmt als die größten gemeinschaftlichen
Teiler von $\partial$ mit $n'$ und $n''$, und darauf wird $c'$ nach
dem Modul $a'$, $c''$ nach dem Modul $a''$ bestimmt aus den beiden
letzten Kongruenzen (4).

Es kommt nun vor allem darauf an, zu zeigen, daß die
Kongruenzen (4) erhalten bleiben, wenn die Transformationen (2)
und (3) nach \S 69, (8) bis (12) durch die beiden linearen Transformationen
\[
(\omega, \omega+1), \quad \left(\omega, -\frac{1}{\omega}\right)
\]
umgeformt werden.

Wir setzen nach den erwähnten Formeln:
\weq{5}{
\begin{aligned}
\begin{pmatrix} a, & 0\\ c, & \partial \end{pmatrix}
\begin{pmatrix} 1, & 0\\ 1, & 1 \end{pmatrix}
&=
\begin{pmatrix} 1, & 0\\ \lambda, & 1 \end{pmatrix}
\begin{pmatrix} a, & 0\\ c_1, & \partial \end{pmatrix}\\
\begin{pmatrix} a', & 0\\ c', & \partial' \end{pmatrix}
\begin{pmatrix} 1, & 0\\ 1, & 1 \end{pmatrix}
&=
\begin{pmatrix} 1, & 0\\ \lambda', & 1 \end{pmatrix}
\begin{pmatrix} a', & 0\\ c_1', & \partial' \end{pmatrix}\\
\begin{pmatrix} a'', & 0\\ c'', & \partial'' \end{pmatrix}
\begin{pmatrix} 1, & 0\\ 1, & 1 \end{pmatrix}
&=
\begin{pmatrix} 1, & 0\\ \lambda'', & 1 \end{pmatrix}
\begin{pmatrix} a'', & 0\\ c_2'', & \partial'' \end{pmatrix}. % [sic] Source prints c_2''; parallelism suggests c_1''.
\end{aligned}
}
\weq{6}{
\begin{aligned}
\begin{pmatrix} a, & 0\\ c, & \partial \end{pmatrix}
\begin{pmatrix} 0, & 1\\ -1, & 0 \end{pmatrix}
&=
\begin{pmatrix} \alpha, & \beta\\ \gamma, & \delta \end{pmatrix}
\begin{pmatrix} a_2, & 0\\ c_2, & \partial_2 \end{pmatrix}\\
\begin{pmatrix} a', & 0\\ c', & \partial' \end{pmatrix}
\begin{pmatrix} 0, & 1\\ -1, & 0 \end{pmatrix}
&=
\begin{pmatrix} \alpha', & \beta'\\ \gamma', & \delta' \end{pmatrix}
\begin{pmatrix} a_2', & 0\\ c_2', & \partial_2' \end{pmatrix}\\
\begin{pmatrix} a'', & 0\\ c'', & \partial'' \end{pmatrix}
\begin{pmatrix} 0, & 1\\ -1, & 0 \end{pmatrix}
&=
\begin{pmatrix} \alpha'', & \beta''\\ \gamma'', & \delta'' \end{pmatrix}
\begin{pmatrix} a_2'', & 0\\ c_2'', & \partial_2'' \end{pmatrix}.
\end{aligned}
}
